1. Beschreibende Fragen:
Diese Fragen zielen darauf ab, einen Zustand oder ein Phänomen zu beschreiben. 

    Beispiel: „Welchen Einfluss hat die Digitalisierung auf das Controlling bei der Gruner Automotive GmbH?“ 

2. Erklärende Fragen:
Sie untersuchen die Ursachen und Zusammenhänge eines Sachverhalts. 

    Beispiel: „Warum wurden im Jahr 2014 im Bundesland Bayern Mindestabstände zwischen Wohnsiedlungen und Windrädern beschlossen?“ 

3. Bewertende Fragen:
Hierbei geht es um die Bewertung, ob etwas gut, sinnvoll, wirksam oder geeignet ist. 

    Beispiel: „Wie erfolgreich war die Einführung des Rechts auf einen Betreuungsplatz für Kinder unter drei Jahren im Hinblick auf das Ziel, den Anteil erwerbstätiger Frauen zu steigern?“ 

4. Prognostizierende Fragen:
Diese Fragen fokussieren auf Vorhersagen über zukünftige Entwicklungen oder Ereignisse. 

    Beispiel: „Wie wird sich der Arbeitskräftebedarf im Bankensektor durch die zunehmende Digitalisierung von Finanzdienstleistungen in den nächsten zehn Jahren entwickeln?“ 


Wer? Bezieht sich auf die Akteure oder Personen.            - Team der Nutzer
Was? Beschreibt das Forschungsobjekt und Ziel.              - Verbesserung der Usability unter Austausch des Systems
Wann? Legt den zeitlichen Rahmen fest.                      - im täglicher Monitoring-Prozess
Warum? Erkundet die Ursachen und Motivationen.              - Modernisierung der Legacy-Plattform
Wie? Erfragt die Methode oder den Prozess.                  - Nutzerstudie